\cotefile{34}
\scanpage{67}
\mspage{II 15}
% archive sheet number: 66
% lines 1–4
\deleted{\unsure{choisissons} $K'$ assez grand, et utilisons le cor. 2 précédent, \unsure{pour} \deleted{\unsure{donc}} \unsure{supposer} que l'action de $\operatorname{Gal}(\overline{K'}/K')$ est \unsure{triviale}. Changeant de notations, on peut supposer alors $K'=K$. Or l'action \unsure{en est triviale}.}

% lines 5–6
Il suffit d'établir le résultat suivant (qui a été signalé par Serre) :

% lines 7–13
\uline{Lemme} Soit $V$ anneau de val. discrète, corps fractions $K$, $A$ schéma abélien sur $K$, \inserted{$\ell$ un nb premier distinct de la car. de \uline{$K$}, (mais pas néc. de la car. du corps résiduel)} supposons que l'action du groupe d'inertie \deleted{\unsure{sur}} $I\subset\operatorname{Gal}(\bar K/K)$ sur $E=T_\ell(A(\bar K))\otimes_{\unsure{\mathbb Z_\ell|Z_\ell}}\unsure{\mathbb Q_\ell|Q_\ell}$ soit \deleted{\unsure{triviale}} \deleted{\unsure{abélienne}} demi-simple, \deleted{\unsure{(dans ce cas abélienne, alors}} alors \deleted{$A$ a bonne} (et alors seulement) $A$ a bonne réduction potentielle''.

% line 13, left margin
% slanted
\margin{\uline{Pour la démonstration} \uline{\unsure{renvoyons|renvoyer} à exposé} \uline{ultérieur}}

% lines 14–17
Le cas de la car. résiduelle $\ne\ell$ est évident, car l'action de $I$ étant un. unipotente (\unsure{J|I}), \inserted{est semi-simple, ssi elle} doit être un. triviale. \deleted{\unsure{et}} Le cas de la car résiduelle $\ell$ (donc $K$ de car zéro) demande un traitement spécial.

% lines 18–24
On peut supposer $V$ \deleted{hensélien, et \uline{complet}, et, après \unsure{ext}} \inserted{à \uline{corps résiduel sép. clos}.} Considérons le modèle de Néron \deleted{\unsure{de}} $G$ de $A$ sur $V$ [\quad]. \deleted{Alors} Le th. de réduction \unsure{$\frac12$ st. R.}, qui sera prouvé ultérieurement, indique que (quitte à faire une extension finie sur $K$) on peut supposer que $G_{\unsure{R|\kappa|k}}^0$ est ext. d'une V.A. par un tore. \inserted{$T$ de dimension $r$} On sait donc que $T_\ell(A(\bar K))$ est extension d'un g.
